\chapter{Development Methodology}

\section{Introduction}

The development methodology defines the approach followed for planning, developing, testing, and refining the Hostel Administration System. Since the requirements of a software system can evolve during development, an iterative methodology was selected instead of a strictly sequential development model.

The Hostel Administration System was developed using an Agile approach with Scrum-inspired development practices. The project was divided into smaller development activities, with each activity focusing on a specific group of features. The functionality was implemented incrementally and tested throughout the development process.

The methodology provided a structured approach for managing the project while also allowing implementation issues and design decisions to be addressed as they were encountered.

\section{Agile Development}

Agile software development is an iterative and incremental approach in which software is developed through repeated cycles of planning, implementation, testing, and refinement. Instead of attempting to implement the complete system in a single development phase, the overall project is divided into smaller functional units.

The Agile approach is particularly suitable for projects where requirements may evolve during development. It allows developers to receive feedback from the implemented system and incorporate improvements into subsequent development cycles.

For the Hostel Administration System, an iterative approach was useful because several components of the system depended on one another. For example, the backend API and database schemas had to be established before the frontend components could be fully integrated, and role-based authentication was required before protected administrative routes could be implemented.

The major principles applied during development were:

\begin{itemize}
    \item Incremental development of system functionality.
    \item Continuous testing during implementation.
    \item Early identification and correction of errors.
    \item Separation of development tasks into manageable units (e.g., backend vs frontend).
    \item Continuous refinement of the system based on implementation results.
\end{itemize}

\section{Scrum Framework}

Scrum is an Agile framework that organises software development into short, iterative development cycles known as sprints. A sprint focuses on implementing a selected set of features or tasks within a defined period.

The Scrum framework generally involves a product backlog, sprint planning, development activities, review, and retrospective evaluation. Although the Hostel Administration System is developed by a smaller team rather than a large commercial software product, Scrum principles provide a useful structure for organising its development.

The major Scrum concepts relevant to the project are:

\begin{itemize}
    \item \textbf{Product Backlog:} A collection of features and tasks required to develop the system.
    
    \item \textbf{Sprint:} A limited development period during which a selected set of backlog items is implemented.
    
    \item \textbf{Sprint Planning:} Identification and prioritisation of tasks to be completed during a development cycle.
    
    \item \textbf{Sprint Review:} Evaluation of the functionality completed during the development cycle.
    
    \item \textbf{Retrospective:} Reflection on problems encountered and possible improvements for subsequent development activities.
\end{itemize}

\section{Product Backlog}

The product backlog for the Hostel Administration System was created by identifying the major functional and technical requirements of the proposed system. These requirements were then divided into smaller implementation tasks.

The initial backlog included the following major items:

\begin{enumerate}
    \item Backend project setup (Node.js, Express) and frontend setup (React, Vite).
    \item MongoDB database configuration and schema design.
    \item Implementation of user registration and JWT-based login.
    \item Backend API development for managing rooms and students.
    \item Frontend authentication state management and protected routing.
    \item Development of the administrative dashboard.
    \item Implementation of the room management interface (Add/Edit rooms).
    \item Implementation of the student allocation interface.
    \item Frontend and backend integration for seamless data flow.
    \item Middleware configuration for role-based security.
    \item Functional testing, debugging, and UI refinement.
\end{enumerate}

The backlog could be refined as development progressed. Tasks that depended on previously implemented functionality were scheduled after the required foundation had been established.

\section{Sprint Planning}

Sprint planning involved selecting a manageable set of tasks from the development backlog. Tasks were arranged according to their dependencies and importance.

The development initially concentrated on establishing the database architecture and backend API foundation. Once the core endpoints and authentication mechanisms were functional, development shifted to the frontend structure. Protected navigation, dashboard functionality, and management forms were subsequently implemented.

The later development stages concentrated on integrating the frontend React components with the Express API and refining the user interface.

The general progression of the development process is illustrated below.

\begin{center}
\begin{tikzpicture}[
    node distance=1.2cm,
    every node/.style={
        draw,
        rounded corners,
        align=center,
        minimum width=5cm,
        minimum height=0.9cm
    },
    arrow/.style={->, thick}
]

\node (planning) {Identify Requirements\\and Create Backlog};

\node (sprint) [below=of planning]
{Select Tasks for Sprint};

\node (development) [below=of sprint]
{Implement Selected Features};

\node (testing) [below=of development]
{Test and Identify Issues};

\node (review) [below=of testing]
{Review and Refine};

\node (next) [below=of review]
{Select Next Development Tasks};

\draw[arrow] (planning) -- (sprint);
\draw[arrow] (sprint) -- (development);
\draw[arrow] (development) -- (testing);
\draw[arrow] (testing) -- (review);
\draw[arrow] (review) -- (next);

\draw[arrow] (next.east) -- ++(2,0) |- (sprint.east);

\end{tikzpicture}
\end{center}

\noindent
Figure 3.1: Iterative Development Process

The iterative nature of this process allowed problems discovered during testing to influence subsequent development tasks.

\section{Sprint Execution}

During sprint execution, the selected tasks were implemented and tested incrementally. Development was performed by first establishing the foundation required by later features.

The initial development activities involved creating the Express backend, configuring MongoDB connections using Mongoose, and defining the data schemas for users, rooms, and students.

The next stage involved implementing backend authentication. JSON Web Tokens (JWT) and password hashing (e.g., bcrypt) were integrated to support secure user registration and login. API endpoints were protected using custom middleware to verify tokens and user roles.

After the backend foundation was established, the React frontend was configured. An authentication context was created to manage the user's login state globally. Protected routes were introduced to ensure that administrative pages could not be accessed by unauthenticated users.

The dashboard was then implemented as the main authenticated interface. Following this, the room and student management pages were developed, connecting frontend forms to the corresponding backend API endpoints.

\section{Development Stages}

The development of the Hostel Administration System can be divided into the following major stages.

\subsection{Stage 1: Environment and Database Setup}

The initial stage established the full-stack project environment. The backend Node.js application and the frontend React/Vite application were initialized. MongoDB was configured, and Mongoose schemas were designed for the core entities (Users, Rooms, Students).

\subsection{Stage 2: Backend API and Authentication}

RESTful API routes were developed using Express. Authentication endpoints (registration, login) were implemented using JWT for secure, stateless session management. Role-based authorisation middleware was created to restrict access to specific administrative actions.

\subsection{Stage 3: Frontend Routing and Auth Context}

The React application structure was defined using React Router. An authentication context was implemented to store the JWT and user details on the client side. Protected navigation components were created to handle redirects for unauthorised access attempts.

\subsection{Stage 4: Core Management Modules}

API endpoints for CRUD (Create, Read, Update, Delete) operations on rooms and students were finalised. On the frontend, interfaces such as the Add Room and Edit Room pages were developed to allow administrators to manage hostel capacity and occupancy. 

\subsection{Stage 5: Integration and Data Fetching}

The frontend components were connected to the backend API using HTTP clients (e.g., Axios or Fetch). React state and effect hooks were utilised to load and display dynamic data, ensuring the UI reflected the current state of the database.

\subsection{Stage 6: Dashboard Development}

A centralised dashboard was developed to provide a consolidated view of hostel operations. The dashboard aggregates data from the API to display key metrics, such as total rooms, available beds, and active students.

\subsection{Stage 7: Testing and Refinement}

Testing was carried out throughout development rather than being restricted to the final stage. API endpoints were tested using tools like Postman, while frontend components, routing, form validation, and database operations were tested as features were implemented.

During development, issues related to CORS (Cross-Origin Resource Sharing), token expiration, and asynchronous data loading were identified and corrected. This iterative testing process helped ensure that the backend and frontend communicated correctly before subsequent functionality was added.

\section{Sprint Review}

At the end of a development cycle, the implemented functionality was reviewed to determine whether the intended features were working correctly. The review involved interacting with the application through its user interface and verifying the expected behaviour against the database state.

For example, after implementing authentication, the login process was tested to ensure the JWT was correctly issued and stored, and that protected API routes correctly rejected requests without a valid token. Similarly, the room management module was reviewed by adding and editing rooms and confirming the changes in the MongoDB database.

The review process helped identify incomplete functionality and integration issues before proceeding with further development.

\section{Retrospective and Refinement}

Retrospective evaluation was used to identify problems encountered during development and determine improvements for subsequent work.

One example involved handling database relationships. Initially, room and student records may have been managed independently, leading to potential inconsistencies when a student was assigned a room. The approach was refined by updating the Mongoose schemas to use object references (ObjectIds) and implementing transaction-like logic in the API controllers to ensure that room occupancy counts were accurately updated when students were allocated.

This demonstrates the usefulness of iterative development. Instead of treating the initial implementation as final, architectural improvements discovered during testing were incorporated into the subsequent development process.

\section{Methodology Applied to the Hostel Administration System}

The development methodology was applied to the project through a sequence of incremental development activities. Each stage produced functionality that served as a foundation for subsequent stages.

The overall development progression can be summarised as:

\begin{enumerate}
    \item Identify the requirements of the hostel administration system.
    \item Establish the backend (Node/Express) and frontend (React/Vite) environments.
    \item Configure the MongoDB database and design data schemas.
    \item Implement JWT authentication and backend security middleware.
    \item Develop RESTful API endpoints for system resources.
    \item Implement frontend authentication context and protected routing.
    \item Develop the core administrative interfaces (rooms, students).
    \item Integrate the frontend UI with the backend API.
    \item Develop the authenticated dashboard and analytics views.
    \item Test and refine the implemented functionality, addressing edge cases and UI improvements.
\end{enumerate}

The project documentation maintained during development records the activities and progress associated with the different stages, providing evidence of the iterative nature of the development process.

\section{Advantages of the Selected Methodology}

The Agile and Scrum-based approach provided several advantages for the development of the Hostel Administration System.

\begin{itemize}
    \item \textbf{Incremental development:} Features could be implemented individually (e.g., API first, then UI) rather than developing the entire application monolithically at once.

    \item \textbf{Early testing:} Errors in database queries or API logic could be detected while the corresponding feature was still being developed.

    \item \textbf{Flexibility:} Development tasks could be adjusted when implementation revealed new requirements, such as adding specific role-based permissions.

    \item \textbf{Manageability:} Dividing a full-stack project into smaller tasks made the development process easier to organise.

    \item \textbf{Continuous refinement:} Feedback from frontend integration testing could be immediately incorporated into backend API adjustments.

    \item \textbf{Risk reduction:} Technical problems, such as cross-origin request issues, could be identified before they affected large portions of the system.
\end{itemize}

\section{Limitations of the Methodology}

Although the Agile approach was highly useful for the project, its implementation within an academic or small-scale project has certain limitations. The development was performed by a small team (or individual) and therefore did not involve the full set of Scrum roles and ceremonies normally found in a professional Scrum environment.

For example, the project does not require a separate Scrum Master, Product Owner, and development team. Planning, development, testing, and review activities are instead performed as part of the individual developer's workflow.

Nevertheless, the underlying principles of incremental full-stack development, continuous testing, task prioritisation, and iterative refinement remain highly applicable and beneficial to the project's success.

\section{Chapter Summary}

This chapter described the development methodology adopted for the Hostel Administration System. An Agile approach incorporating Scrum principles was used to organise the full-stack development process into manageable and iterative activities.

The development progressed from backend environment setup and MongoDB schema design to REST API development, JWT authentication, frontend React integration, protected navigation, and administrative dashboard development. Testing and refinement were carried out continuously across both the frontend and backend throughout the development process.

The methodology provided a structured yet flexible approach for developing the system and established a clear connection between the project's development activities and the functionality implemented in the final MERN-stack application.
